\documentclass[11pt]{article}
\input{preamble}
\input{macros}
\addbibresource{refs.bib}
\linespread{1.11}

\title{Categorical Message Passing Language (CaMPL)}
\author{
    Robin Cockett\thanks{University of Calgary} \and Daniel Hashimoto\thanks{Universidade Federal do Rio de Janeiro} \and Alexanna Little* \and Melika Norouzbeygi* \and Priyaa Varshinee Srinivasan\thanks{Tallinn University of Technology}  \and \ldots}
%\def\authorrunning{D.\ I.\ Spivak \& P.\ V.\ Srinivasan}
\date{\today}

\begin{document}

\maketitle

\begin{abstract}
This document is a simple introduction to  a novel concurrent programming language called the 
categorical message passing language built on the mathematics of linear actegories. 
Categorical Message Passing Language, to the best of our knowledge, 
the only\footnote{Is this true?} concurrent programming language with mathematical underpinnings. 
It is a functional-style programming language with the abstract machine built on Haskell. 
The current version of CaMPL compiler is to be considered a proto-alpha, a proof-of-concept 
of a programming language built from the mathematics of linear actegories.
\end{abstract}

%\addcontentsline{toc}{section}{References}

\tableofcontents

\section{Introduction}

Explain current state of concurrency and why CaMPL is novel (Curry-Howard-Lambek-like correspondence), and its history.

\note{Xanna note:
i am assuming introduction will establish justification for why we have submitted this paper to this venue.
and that we will mention things like there is a beautiful type system and that we are doing message passing and that there is a categorical semantics}

As is tradition, we will introduce CaMPL with a ``Hello World!'' example program. The ``Hello World!'' example demonstrates communication between a process and the outside world using the {\tt Console} channel type. 

The following code snippet shows a CaMPL process named {\tt helloworld} printing {\tt ``Hello World!''} by sending the string as a message on a channel named {\tt console}. 
The {\tt console} channel has the type {\tt Console} which is a special channel type hard-coded in the compiler to connect the program to the terminal from which it is run.
\begin{lstlisting}
proc helloworld :: | Console => =
	| console => -> do
		hput ConsolePut on console
		put "Hello World!" on console    -- sends message to the console
		hput ConsoleClose on console
		halt console                     -- closes console channel and halts
\end{lstlisting}
The \tc{process} \lstinline$helloworld$ communicates to an in-built printing process via a channel of type \lstinline*Console*. In CaMPL, concurrent types are \tc{(co)protocols}, the instantiations of which are called as \tc{channels}. The \tc{process command} \lstinline$hput$ may seem strange at the first sight --- read as `handle put',  it initiates a channel for a  communication session. \tc{Protocol handles} are analogous to session types. As indicated by their names, the \lstinline$put$ process command sends a message, and the \lstinline$halt$ process command halts the calling process.

The {\tt helloworld} process can be invoked by calling it in the main process, {\tt run}, and giving it the {\tt console} channel as follows:
\begin{lstlisting}
 proc run = 
	| console => -> 
		helloworld( | console => )  -- creates helloworld process
\end{lstlisting}

The channel to the console can only be created by the main process.

\section{Designing a simple database server}
% Describe in-built channel types - Put, Get, TopBot
% Describe process command - plug
% Describe channel polarities and their roles
% A drop of mathematics and logic

As we become more familiar with CaMPL, let us begin by designing a simple database server in which a client and a server will exchange one message each. 

In order to arrange such a communication to proceed, we need to connect the client and the server, which is facilitated by the \lstinline$plug$ process command. 
\begin{lstlisting}
proc run = 
	| => -> 
		plug 
			client( | => ch )
			server( | ch => )
\end{lstlisting}

The above {\tt run} process, in Line 3, creates a {\tt client} and a {\tt server} process, and plugs them together via a named {\tt ch}. The CaMPL compiler has type-inference ability and will resolve the type of the channel {\tt ch}. 

It may be noticed that, in line 4, that {\tt ch} is placed after "\lstinline$=>$" for the client, and in line 5, that, {\tt ch} is placed before "\lstinline$=>$" for the server. Whether {\tt ch} occurs before or after the arrow determines its \tc{channel polarity} --- the different polarities signifies different ends of the channel. In this case, {\tt ch} is of \tc{output polarity}, also written as \ttc{(+)} for the {\tt client}. Dually, {\tt ch} is of \tc{input polarity}, also written as \ttc{(-)} for the {\tt server}. 

Note: The choice of terminology -- input and output -- polarities is best understood by imagining a process as a box, with wires entering it and wires leaving it (with respect to the direction of gravity). We call incoming wires as input polarity, and outgoing as output polarity. 

Thus, it must be obvious now that CaMPL allows processes to be plugged together only along channels of same type and opposite polarities. 

Let us now to code the client process:
\begin{lstlisting}
proc client :: | => Put(Int | Get( [Char] | TopBot)) =
	| => ch -> do
		-- Do something --
		put val on ch    -- send a message
		get msg on ch    -- receive a message 
		halt ch          -- close channel and halt process
\end{lstlisting}
% Explain type versus polarity.
\lstinline$Put$, \lstinline$Get$, and \lstinline$TopBot$ are in-built concurrent types for channels. The \lstinline$Put$ and \lstinline$Get$ types is used to exchange messages, while the \lstinline$TopBot$ type is used to end communication.
\begin{itemize}
\item Line 1 of \lstinline$proc$ {\tt client} clarifies that the {\tt client} process has an output polarity channel of compound type \lstinline$Put(Int | Get( [Char] | TopBot))$. 
\item Line 2 assigns variable name {\tt ch} to the output polarity channel. 
\item Line 4 sends the integer value of {\tt val} over {\tt ch} which is of type \lstinline$Put(Int | Get( [Char] | TopBot))$. On completion, the {\tt ch} proceeds with type \lstinline$Get( [Char] | TopBot)$.
\item Line 5 receives a string over {\tt ch} which is now of type \lstinline$Get( [Char] | TopBot)$, and binds the string value to the variable {\tt msg}. On completion, the {\tt ch} proceeds with type \lstinline$TopBot)$.
\item Line 6 closes {\tt ch} which is now of type {\tt TopBot} (hence, no more communication over {\tt ch}) and halts the {\tt client} process.
\end{itemize}

It will become clear in a moment the role played by polarities in communication, when we code the server process. But, for now, it must be clear that a process may both send and receive over an output polarity channel.

Let us design the server process:
\begin{lstlisting}
proc server :: | Put(Int | Get( [Char] | TopBot)) => =
	| ch => -> do
		get inval on ch     -- receive message
		-- Do something with the input  
		put outmsg on ch    -- send message 
		halt ch             -- close channel and halt process
\end{lstlisting}
\begin{itemize}
\item Line 1 gives the type signature of \lstinline$proc$ {\tt server}. The {\tt server} process has an input polarity channel of compound type \lstinline$Put(Int | Get( [Char] | TopBot))$. 
\item Line 2 assigns variable name {\tt ch} to the input polarity channel. 
\item Line 4 receives an integer value over channel {\tt ch} which is of type \lstinline$Put(Int | Get( [Char] | TopBot))$, and binds it to the variable {\tt inval}. On completion, the {\tt ch} proceeds with type \lstinline$Get( [Char] | TopBot)$.
\item Line 5 sends a string over channel {\tt ch} which is now of type \lstinline$Get( [Char] | TopBot)$. On completion, the {\tt ch} proceeds with type \lstinline$TopBot)$. 
\end{itemize}

Note that, when the process command `\lstinline$put val on ch$' of the {\tt client} process (Line 4) is matched by the process command `\lstinline$get inval on ch$' in the {\tt server} process. This command pair allows communication of a single message. We call the pair {\tt (put, get)} as a \tc{complementary pair} -- meaning the commands are used in pairs one at each end of the same channel. This arises a question, given a complementary pair, which command should be used at which end of the channel? This is where polarities come into play.

Given a channel of type \lstinline$Put$, the process which has the {\bf output polarity} end may \lstinline$put$ a message on the channel, and correspondingly, the process which has the {\bf input polarity} end may \lstinline$get$ a message over the channel, and not vice versa. Given a channel of type \lstinline$Get$, the process which has the {\bf input polarity} end may \lstinline$put$ a message on the channel, and correspondingly, the process which has the {\bf output polarity} end may \lstinline$get$ a message over the channel, and not vice versa. Thus, given either a \lstinline$Put$ or a \lstinline$Get$ type, the polarities allow unambiguous distinction of the role of a process (if it is a sender or a receiver). 

\note{@Xanna, I think you are best suited to write the following paragraph :)}
The notion polarities arise from the underlying mathematics of linear actegories. The linear actegories are equipped with an adjunction parametrized by the sequential datatypes. The following four sequent rules encode the usage of {\tt(put,get)} process commands over \lstinline$Put$ and \lstinline$Get$ types.

\section{Concurrent type system: Handling multiple sequential queries}

\section{Split and Fork: Asynchronous querying and response}

% This is a good place to talk about deadlocks.

\section{Races: Handling multiple concurrent queries}

\section{Higher-order message passing: Dyanmic database}

\section{Mathematical underpinning}

\section{Discussion and future work}


\begin{appendix}
\section{Examples: Complete programs}
\end{appendix}

\end{document}
